\documentclass[11pt]{article}

\usepackage[T1]{fontenc}
\usepackage{amsmath,amssymb,amsthm}
\usepackage[margin=1in]{geometry}
\usepackage[colorlinks=true,linkcolor=blue,citecolor=blue,urlcolor=blue]{hyperref}

\newtheorem{theorem}{Theorem}
\newtheorem{proposition}[theorem]{Proposition}
\newtheorem{lemma}[theorem]{Lemma}
\newtheorem{conjecture}[theorem]{Conjecture}
\theoremstyle{remark}
\newtheorem{remark}[theorem]{Remark}

\newcommand{\R}{\mathbb R}
\newcommand{\Pj}{\mathbb P}
\newcommand{\ip}[2]{\left\langle #1,#2\right\rangle}
\newcommand{\norm}[1]{\left\|#1\right\|}
\newcommand{\supp}{\operatorname{supp}}
\newcommand{\weakto}{\rightharpoonup}
\setlength{\emergencystretch}{2em}

\title{The q4 Terminal Defect Forces a Zero-Trace Oseen Entrance State}
\author{John Lawson and Codex}
\date{25 July 2026}

\begin{document}
\maketitle

\begin{abstract}
This round bundles the entire conditional q4 terminal energy defect into
one same-trajectory projected-Oseen adjoint.  Finite-band terminal data
\(g_j\weakto0\) are extracted from the forced ultraviolet energy floor.
Their backward adjoints converge to
\[
 a\in L^\infty_tL^2_\sigma\cap L^2_t\dot H^1_\sigma
\]
with constant nonzero primal pairing
\[
 \ip{u(t)}{a(t)}=c_0>0,
\]
although \(a(t)\weakto0\) as \(t\uparrow T^*\).  Thus \(a\) is a
reverse-time entrance state with zero weak trace and a fixed strong-trace
defect; finite kinetic energy alone proves the weak trace.  The associated
terminal cross-defect measure is nonzero and is supported on the same
singular slice as the primal energy defect.

For the bounded-band subcase, a literal piecewise adjoint has a fixed
\(L^2\) reset cost at every event, so its source variation grows linearly.
Pure heat is negligible at the clock band.  The elementary low-pass strain
estimate controls only frequencies
\(\mu_j\asymp m_j^{9/4}j^{1/2}\), below the clock
\(\Lambda_j\asymp m_j^{7/3}j^{2/3}\) by
\(m_j^{1/12}j^{1/6}\).  The remaining gate is a new terminal
strong-trace theorem for the pressure-coupled, self-generated Oseen
system.  No Clay alternative is proved.
\end{abstract}

\section{Epistemic status and inputs}

The statements proved here are conditional repository theorems, subject
to external mathematical review.  Their antecedent is the
energy-efficient exact q4 branch of the repository's Type-II reduction.
That branch is not known to follow from arbitrary smooth Clay data.

The imported Navier--Stokes partial-regularity statement is the classical
Caffarelli--Kohn--Nirenberg theorem~\cite{CKN82}.  The scalar
divergence-free-drift comparison is from Qian and Xi~\cite{QX19}.
Their peer-reviewed theorem is contextual only; the projected vector
adjoint below is not their scalar equation.

The Clay problem remains unsolved.

\section{Terminal detector}

Let \(u\) be smooth on \([0,T)\), with a Leray--Hopf continuation through
the putative first singular time \(T\).  Put
\[
 u_*:=u(T),
 \qquad
 w_j:=u(s_j)-u_*,
 \qquad
 U_2:=\sup_{0<t<T}\norm{u(t)}_2.
\]
The preceding critical-clock theorem gives
\begin{equation}
 \label{eq:highpass}
 \norm{\Pi_{>\Lambda_j}w_j}_2\geq d_1>0,
 \qquad
 \Lambda_j\longrightarrow\infty.
\end{equation}

\begin{lemma}[Finite-band terminal detector]
\label{lem:detector}
There are finite \(\Omega_j>\Lambda_j\) and solenoidal data
\[
 g_j
 :=
 \Pi_{\Lambda_j<|\xi|\leq\Omega_j}w_j
\]
such that
\begin{equation}
 \label{eq:g-floor}
 \norm{g_j}_2\geq\frac{d_1}{2},
 \qquad
 g_j\weakto0
 \quad\hbox{in }L^2.
\end{equation}
After extraction,
\begin{equation}
 \label{eq:terminal-pair}
 \ip{u(s_j)}{g_j}\longrightarrow c_0,
 \qquad
 c_0\geq\frac{d_1^2}{4}>0.
\end{equation}
\end{lemma}

\begin{proof}
For fixed \(j\), monotone convergence of the Fourier tail permits a
finite \(\Omega_j\) for which \eqref{eq:g-floor} holds.  The lower
Fourier support tends to infinity, so \(g_j\weakto0\).
Furthermore, \(\norm{g_j}_2\leq\norm{w_j}_2\leq2U_2\).

Orthogonality of the projection gives
\[
 \ip{u(s_j)}{g_j}
 =
 \norm{g_j}_2^2+\ip{u_*}{g_j}.
\]
Moreover,
\[
 |\ip{u_*}{g_j}|
 \leq
 \norm{\Pi_{>\Lambda_j}u_*}_2\norm{g_j}_2
 \longrightarrow0.
\]
The pairings are bounded; select a convergent subsequence and use
\eqref{eq:g-floor}.
\end{proof}

\begin{remark}
The individual upper cutoffs may diverge arbitrarily quickly.  Therefore
Lemma~\ref{lem:detector} uses only the high-pass floor and applies before
choosing between bounded clock-band mass and super-clock escape.
\end{remark}

\section{Projected backward Oseen identities}

For each \(j\), let \(a_j\) solve
\begin{equation}
 \label{eq:adjoint}
 \partial_ta_j+\nu\Delta a_j
 +\Pj\bigl((u\cdot\nabla)a_j\bigr)=0,
 \qquad
 \nabla\cdot a_j=0,
 \qquad
 a_j(s_j)=g_j.
\end{equation}
This is a stable parabolic problem in reverse time.  In pressure form,
\begin{equation}
 \label{eq:adjoint-pressure}
 \partial_ta_j+\nu\Delta a_j
 +(u\cdot\nabla)a_j+\nabla q_j=0,
 \qquad
 -\Delta q_j
 =
 \partial_i u_\ell\,\partial_\ell(a_j)_i.
\end{equation}
Here ``projected Oseen adjoint'' means the adjoint of the frozen-drift
operator
\[
 \mathcal L_u v
 :=
 \partial_t v-\nu\Delta v+\Pj((u\cdot\nabla)v).
\]
It is not the adjoint of the full Fr\'echet linearisation of
Navier--Stokes, which also varies the advecting field.

\begin{lemma}[Energy and exact primal pairing]
\label{lem:identities}
For \(0\leq t\leq s_j\),
\begin{equation}
 \label{eq:adjoint-energy}
 \boxed{
 \norm{a_j(t)}_2^2
 +2\nu\int_t^{s_j}\norm{\nabla a_j(\tau)}_2^2\,d\tau
 =
 \norm{g_j}_2^2
 }
\end{equation}
and
\begin{equation}
 \label{eq:pairing}
 \boxed{
 \ip{u(t)}{a_j(t)}
 =
 \ip{u(s_j)}{g_j}.
 }
\end{equation}
\end{lemma}

\begin{proof}
Pair \eqref{eq:adjoint-pressure} with \(a_j\).  The drift is skew
because \(u\) is divergence free, and the pressure is orthogonal to
\(a_j\).  Integration backwards from \(s_j\) gives
\eqref{eq:adjoint-energy}.

The projected Navier--Stokes equation is
\[
 \partial_tu
 =
 \nu\Delta u-\Pj((u\cdot\nabla)u).
\]
Differentiate \(\ip{u}{a_j}\).  The two diffusion terms cancel by
self-adjointness, while
\[
 \ip{(u\cdot\nabla)u}{a_j}
 =
 -\ip{u}{(u\cdot\nabla)a_j}.
\]
All pressure pairings vanish.  Hence the derivative is zero.
\end{proof}

\section{One zero-trace entrance state}

\begin{theorem}[Projected-Oseen entrance]
\label{thm:entrance}
After passing to a subsequence, there exists
\[
 a\in
 L^\infty(0,T;L^2_\sigma(\R^3))
 \cap
 L^2(0,T;\dot H^1_\sigma(\R^3))
\]
which solves
\begin{equation}
 \label{eq:limit-adjoint}
 \partial_ta+\nu\Delta a
 +\Pj((u\cdot\nabla)a)=0
\end{equation}
distributionally on \((0,T)\), and satisfies
\begin{equation}
 \label{eq:constant-limit-pair}
 \boxed{
 \ip{u(t)}{a(t)}=c_0>0
 \quad(0<t<T).
 }
\end{equation}
Nevertheless,
\begin{equation}
 \label{eq:weak-zero}
 \boxed{
 a(t)\weakto0
 \quad\hbox{in }L^2
 \quad(t\uparrow T).
 }
\end{equation}
In particular,
\begin{equation}
 \label{eq:norm-floor}
 \boxed{
 \inf_{0<t<T}\norm{a(t)}_2
 \geq\frac{c_0}{U_2}>0.
 }
\end{equation}
\end{theorem}

\begin{proof}
The energy identity \eqref{eq:adjoint-energy} gives uniform bounds in
\(L^\infty L^2\cap L^2\dot H^1\).  On
\([0,T-\varepsilon]\), the coefficient \(u\) is smooth and the equation
bounds \(\partial_ta_j\) in a negative Sobolev space.  Weak compactness,
weak time continuity, and a diagonal extraction yield \(a\) and permit
passage to \eqref{eq:limit-adjoint}.  Equations
\eqref{eq:terminal-pair} and \eqref{eq:pairing} give
\eqref{eq:constant-limit-pair}.

Let \(h\) be a solenoidal Schwartz field.  Pair
\eqref{eq:adjoint-pressure} with \(h\), integrate from \(t\) to \(s_j\),
and integrate the drift by parts.  The energy bounds for both fields give
\begin{align}
 \label{eq:weak-trace-est}
 |\ip{a_j(t)-g_j}{h}|
 \leq{}&
 C U_2\nu(s_j-t)\norm{\Delta h}_2
 \nonumber\\
 &+
 C U_2^2(s_j-t)\norm{\nabla h}_\infty.
\end{align}
Indeed,
\[
 \norm{a_j}_2\leq2U_2,
 \qquad
 \norm{u\cdot\nabla h}_2
 \leq U_2\norm{\nabla h}_\infty.
\]

Letting \(j\to\infty\), using \(g_j\weakto0\), yields
\[
 |\ip{a(t)}{h}|
 \leq
 C_h(T-t).
\]
This tends to zero as \(t\uparrow T\).  Density and the uniform \(L^2\)
bound prove \eqref{eq:weak-zero}.

Finally, \eqref{eq:constant-limit-pair} and Cauchy--Schwarz give
\[
 c_0
 \leq
 \norm{u(t)}_2\norm{a(t)}_2
 \leq
 U_2\norm{a(t)}_2,
\]
which is \eqref{eq:norm-floor}.
\end{proof}

\begin{remark}[Correct reverse-time interpretation]
Writing \(\tau=T-t\) turns \eqref{eq:limit-adjoint} into a forward
parabolic system.  Its weak initial trace is zero, but its interior
\(L^2\) norm has the floor \eqref{eq:norm-floor}.  The theorem does not
assert that the limit satisfies the initial energy inequality at
\(\tau=0\).  Such an assertion would already give strong attainment of
the zero datum and would rule the state out.
\end{remark}

\section{Cross defect on the singular slice}

Let
\[
 \sigma
 :=
 \{x:(x,T)\text{ is singular for the continuation}\}.
\]
The preceding round proved, conditionally,
\[
 \frac35\leq\dim_{\rm H}\sigma\leq1,
 \qquad
 \mathcal H^1(\sigma)=0.
\]

\begin{proposition}[Terminal cross-defect measure]
\label{prop:cross}
After extraction along the record times, the signed measures
\[
 u(s_j)\cdot a(s_j)\,dx
\]
converge weak-star to a finite signed measure \(\zeta\) with
\begin{equation}
 \label{eq:cross}
 \boxed{
 \zeta(\R^3)=c_0,
 \qquad
 \supp\zeta\subset\sigma.
 }
\end{equation}
If \(\mathcal A\) is a simultaneous local weak-star limit of
\(|a(s_j)|^2dx\), then
\begin{equation}
 \label{eq:a-charges}
 \mathcal A(\sigma)>0.
\end{equation}
\end{proposition}

\begin{proof}
The cross measures have uniformly bounded total variation.  The q4
terminal sequence is uniformly spatially tight in \(L^2\), so
Cauchy--Schwarz makes the cross measures tight.  Their total mass is
\(c_0\) by \eqref{eq:constant-limit-pair}.

If \(K\Subset\R^3\setminus\sigma\), local regularity gives
\[
 u(s_j)\longrightarrow u_*
 \quad\hbox{strongly in }L^2(K),
\]
whereas \eqref{eq:weak-zero} gives \(a(s_j)\weakto0\).  A compactly
supported cutoff near \(K\) therefore makes the cross pairing tend to
zero.  Exhausting the regular set proves the support statement.

Simultaneous measure extraction and Cauchy--Schwarz give
\[
 |\zeta|(B)^2
 \leq
 \mathcal E(B)\mathcal A(B)
\]
for Borel \(B\), where \(\mathcal E\) is the primal terminal energy
measure.  Since \(\zeta\neq0\) and is supported on \(\sigma\),
\eqref{eq:a-charges} follows.
\end{proof}

\section{The fixed reset paid by a piecewise adjoint}

Now restrict to the bounded clock-band alternative.  Its event detectors
\(g_k\) and separated times satisfy
\begin{equation}
 \label{eq:event}
 \ip{u(s_{k+1})-u(s_k)}{g_k}
 \leq-\frac{3\eta}{4}.
\end{equation}
Let \(b_k\) solve the projected backward adjoint on
\([s_k,s_{k+1}]\) with \(b_k(s_{k+1})=g_k\), and put
\[
 r_k:=g_k-b_k(s_k).
\]

\begin{proposition}[Reset floor]
\label{prop:reset}
For every event,
\begin{equation}
 \label{eq:reset-floor}
 \boxed{
 \ip{u(s_k)}{r_k}\geq\frac{3\eta}{4},
 \qquad
 \norm{r_k}_2\geq\frac{3\eta}{4U_2}.
 }
\end{equation}
The canonical forced adjoint which resets at the first \(N\) events has
\(L^2\)-valued source variation at least
\begin{equation}
 \label{eq:variation}
 \sum_{k=1}^N\norm{r_k}_2
 \geq
 \frac{3\eta}{4U_2}N.
\end{equation}
\end{proposition}

\begin{proof}
Lemma~\ref{lem:identities} on the \(k\)-th interval gives
\[
 \ip{u(s_k)}{b_k(s_k)}
 =
 \ip{u(s_{k+1})}{g_k}.
\]
Subtract from \(\ip{u(s_k)}{g_k}\) and use \eqref{eq:event}.
Cauchy--Schwarz gives the norm floor.  Each reset is the impulse
\(r_k\delta_{s_k}\) in the forced equation, and vector-measure total
variation sums the norms.
\end{proof}

\begin{remark}
Proposition~\ref{prop:reset} excludes finite \(L^2\)-variation gluing.
It does not exclude every weighted or cancellation-sensitive square
function.  The unforced limit from Theorem~\ref{thm:entrance} avoids the
jumps by carrying the terminal entrance defect instead.
\end{remark}

\section{Clock-scale cost audit}

With
\[
 M(t):=\norm{u(t)}_{L^{3,\infty}},
\]
the exact record tail gives
\begin{equation}
 \label{eq:tails}
 \int_{s_j}^TM(t)\,dt
 \lesssim
 \frac{m_j^{-9/2}}{j},
 \qquad
 T-s_j
 \lesssim
 \frac{m_j^{-11/2}}{j},
\end{equation}
and
\begin{equation}
 \label{eq:clock}
 \Lambda_j\asymp m_j^{7/3}j^{2/3}.
\end{equation}

\subsection{Heat cannot make the bounded-band reset}

In the bounded-band alternative, \(g_j\) is supported below
\(A\Lambda_j\).  For the backward heat evolution \(H_j\) with terminal
value \(g_j\),
\begin{align}
 \norm{H_j(s_j)-g_j}_2
 &\leq
 \nu(s_{j+1}-s_j)(A\Lambda_j)^2\norm{g_j}_2
 \nonumber\\
 &\lesssim
 A^2m_j^{-5/6}j^{1/3}
 \longrightarrow0.
 \label{eq:heat-small}
\end{align}
Thus pure diffusion of the original clock band cannot account for
\eqref{eq:reset-floor}.

\subsection{The low-pass strain ceiling}

Lorentz Bernstein gives
\[
 \norm{\nabla\Pi_{\leq\mu}u(t)}_\infty
 \lesssim
 \mu^2M(t).
\]
Therefore
\begin{equation}
 \label{eq:strain}
 \int_{s_j}^T
 \norm{\nabla\Pi_{\leq\mu}u(t)}_\infty\,dt
 \lesssim
 \mu^2\frac{m_j^{-9/2}}{j}.
\end{equation}
The largest frequency for which this upper estimate is order one is
\[
 \mu_j\asymp m_j^{9/4}j^{1/2}.
\]
Comparison with \eqref{eq:clock} gives
\begin{equation}
 \label{eq:strain-gap}
 \boxed{
 \frac{\Lambda_j}{\mu_j}
 \asymp
 m_j^{1/12}j^{1/6}
 \longrightarrow\infty.
 }
\end{equation}
At the clock itself, the right side of \eqref{eq:strain} is bounded only
by
\[
 C m_j^{1/6}j^{1/3},
\]
which is not perturbatively small.  This is an upper estimate and not a
claim that the actual strain is large.

\subsection{Pressure}

The pressure \(q_j\) in \eqref{eq:adjoint-pressure} does no global
\(L^2\) work and leaves both exact identities unchanged.  Variable
transport already mixes frequencies; pressure adds nonlocal component
coupling through the enforcement of solenoidality.  A componentwise
scalar maximum principle cannot simply be applied after deleting it.

\section{Comparison with scalar drift theory}

The q4 drift belongs to strong sub-endpoint spaces, but their parabolic
indices remain strictly supercritical.  For \(0<\theta<1\), set
\[
 \frac1{q_\theta}
 =
 \frac{1-\theta}{2}+\frac{\theta}{3}.
\]
Interpolation gives
\[
 \norm{u(t)}_{q_\theta}
 \lesssim
 U_2^{1-\theta}M(t)^\theta,
\]
so
\begin{equation}
 \label{eq:strong-drift}
 u\in L^\ell_tL^{q_\theta}_x
 \qquad
 \text{whenever }\theta\ell<\frac{11}{2}.
\end{equation}
At the limiting time exponent,
\begin{equation}
 \label{eq:scalar-index}
 \frac2\ell+\frac3{q_\theta}
 =
 \frac32-\frac{3\theta}{22}
 >
 \frac{15}{11}
 >1.
\end{equation}

Qian and Xi prove an Aronson-type upper bound for smooth
divergence-free scalar drifts with index in \([1,2)\)
under their stated assumptions~\cite{QX19}.  Their source explicitly
leaves regularity open in the supercritical case and proves
energy-class uniqueness in the critical \(L^\infty_tL^n_x\) case.
Equations \eqref{eq:strong-drift}--\eqref{eq:scalar-index} do not enter
that critical uniqueness theorem.  The projected vector pressure is an
additional distinction.

\section{What changed in this round}

\subsection*{Formal findings, subject to external review}

\begin{enumerate}
\item The entire q4 high-pass defect admits finite-band terminal
detectors with fixed \(L^2\) mass and weak limit zero.
\item Their projected backward adjoints converge to one nonzero
energy-class entrance field whose zero weak terminal trace follows from
finite kinetic energy alone.
\item The constant primal--adjoint pairing produces a nonzero terminal
cross-defect measure supported on the singular slice.
\item Literal piecewise-adjoint sewing pays a fixed \(L^2\) reset at
every bounded-band event.
\item Heat is negligible at the bounded clock band, while the standard
low-pass strain estimate loses the factor
\(m_j^{1/12}j^{1/6}\).
\item The available scalar drift hull approaches parabolic index
\(15/11\), strictly above the critical line.
\end{enumerate}

\subsection*{Closed shortcuts}

\begin{enumerate}
\item Finite \(L^2\)-variation jump gluing cannot bundle infinitely many
event detectors.
\item Pure heat evolution cannot generate the required bounded-band
reset.
\item A perturbative low-pass Lipschitz argument does not reach the
clock frequency.
\item Scalar critical-drift uniqueness cannot be imported at the q4
exponents, and deleting the projected pressure changes the equation.
\end{enumerate}

\subsection*{Open obligations}

\begin{enumerate}
\item Prove terminal strong-trace rigidity for the self-generated
projected Oseen adjoint.
\item Alternatively, exclude a nonzero cross-defect measure supported
on the Caffarelli--Kohn--Nirenberg terminal slice.
\item Find a pressure-sensitive compactness or microlocal budget for the
entrance state.
\item Control divergent normalised energy and clocks at or below the
five-power barrier.
\item Prove one complete Clay alternative for arbitrary admissible data.
\end{enumerate}

\begin{conjecture}[Self-generated Oseen terminal rigidity]
Let \(u\) be a smooth finite-energy Navier--Stokes trajectory before its
first possible singular time and suppose the exact q4 record clock holds.
If
\[
 a\in L^\infty_tL^2_\sigma\cap L^2_t\dot H^1_\sigma
\]
solves \eqref{eq:limit-adjoint} and
\(a(t)\weakto0\) as \(t\uparrow T\), then
\[
 \norm{a(t)}_2\longrightarrow0.
\]
\end{conjecture}

Theorem~\ref{thm:entrance} shows that this conjecture excludes the q4
terminal energy defect.  It is not proved.

\section*{Source provenance}

The pinned arXiv v3 archive for Qian--Xi has SHA-256
\[
\texttt{cf19b869a03ef7554f2675b5dd3e6e9aa8f5041dc39ce4179199ef426e55f50d}.
\]
Its TeX member
\texttt{Divergence\_free\_Qian-Xi-2017-12-15.tex} has SHA-256
\[
\texttt{5e590c76c8e381a0a7b669f9a05ac4179a4c06ce7c9d924c0b2c3b92df082331}.
\]
The supercritical upper-bound theorem is at source lines 217--253; the
explicit supercritical regularity frontier and the critical-only
uniqueness application are at lines 318--335 and 1282--1390.

\begin{thebibliography}{9}
\raggedright

\bibitem{CKN82}
L.~Caffarelli, R.~Kohn, and L.~Nirenberg,
\emph{Partial regularity of suitable weak solutions of the
Navier--Stokes equations},
Communications on Pure and Applied Mathematics 35 (1982), 771--831.

\bibitem{QX19}
Z.~Qian and G.~Xi,
\emph{Parabolic equations with divergence-free drift in space
\(L_t^\ell L_x^q\)},
Indiana University Mathematics Journal 68 (2019), 761--797,
\href{https://doi.org/10.1512/iumj.2019.68.7685}
{doi:10.1512/iumj.2019.68.7685}.

\end{thebibliography}

\end{document}
